\subsection{Integral multiplicities in the Morita Gram identity}
\label{sec:integral-morita-rotation}

The full Morita action Gram identity of
Theorem~\ref{thm:full-morita-gram} has an elementary integral refinement.
The unknown attached to a possible right-action matrix counts simple
objects of the dual category; it is an integer, even when distinct simple
objects have the same action matrix. Retaining this integrality detects
obstructions invisible in the rational nonnegative relaxation.

\begin{proposition}[Morita Gram congruence obstruction]
\label{prop:morita-gram-congruence}
Fix one of the exhaustively enumerated module actions for a forced split
involution algebra $\mathbf1\oplus g$. Let $Ax=b$ be its full Morita Gram
system, including the prescribed first-column multiplicities, with all
possible right-action matrices retained. If there exist $q>1$ and an
integral row vector $y$ such that
\[
 yA\equiv0\pmod q,
 \qquad yb\not\equiv0\pmod q,
\]
then this module action cannot occur in a complex fusion category.
\end{proposition}
\begin{proof}
Every categorification would give a vector $x$ of nonnegative integral
multiplicities satisfying $Ax=b$. Multiplication by $y$ would give
$yb=yAx\equiv0\pmod q$, a contradiction. All entries of $A$ and $b$ come
from the integral fusion and module actions. The argument therefore
requires no pivotal or spherical structure and no positivity assumption
on categorical dimensions.
\end{proof}

For discovery, Gaussian elimination over $\mathbb F_2$ produces the row
vector $y$. Verification only multiplies integer vectors and checks the
two congruences. Repeated equality rows may be removed by a deterministic
rule. The completeness of the right-action candidate list is exactly the
finite Gram enumeration already used in
Section~\ref{sec:full-morita-gram}; no new search cutoff is introduced.

\begin{proposition}\label{prop:round3-higher-stabilizer}
The rings \texttt{R6\_5\_638} and \texttt{R6\_7\_1225} admit no
categorification by a fusion category over $\mathbb C$.
\end{proposition}
\begin{proof}
The unrestricted signed reduction is verified separately for each full
fusion tensor. There are respectively $11\,800$ and $46\,856$ signed
tables, partitioned into $418$ and $1\,560$ sign orbits. Each has one
original gauge orbit. Every other orbit is excluded by its exact
character obstructions; within the original orbit the only admissible
nonzero dimension characters lie in the Frobenius--Perron Galois orbit.
It therefore suffices to exclude pseudo-unitary categorification.

In each case the involution-module enumeration has sixteen branches.
Two fail the full Morita rank bound. Twelve fail
Proposition~\ref{prop:morita-gram-congruence}, with $q=2$. Both remaining
branches force negative stabilizer commutators on the two fixed simple
objects, indexed $4$ and $5$ in the stored tensor. Their cubic rotation
spaces have common dimension $n=4$ for \texttt{R6\_5\_638} and $n=6$
for \texttt{R6\_7\_1225}. The higher-dimensional Clifford rotation
criterion of Section~\ref{sec:stabilizer-rotation} applies. The exponent
prime supports are respectively $\{2,5,13\}$ and $\{2,5,29\}$, so the
third indicators belong respectively to
\[
 \{-2,1\}\qquad\text{and}\qquad\{-3,0\}.
\]
The actual separating certificates already hold on the weaker intervals
$[-2,1]$ and $[-3,3/2]$; they do not rely on selecting a discrete indicator
value.

For either tensor the complete low-induction enumeration gives $91$
frames, of which $31$ have a nonempty character-labelling set. The $50$
labellings and four second-indicator choices yield $200$ exact
height-six systems. Every system has a verified integral separating
functional. The two dimension fields are $\mathbb Q(\sqrt{65})$ and
$\mathbb Q(\sqrt{145})$; their low-column counts are respectively
$(217,225,225,61,61)$ and $(469,481,481,127,127)$. All possible twist
orders are covered by the inherited relative-trace construction.
Thus the two remaining module branches also fail.
\end{proof}

The exact verifier is
\filename{software/round3_integral_morita.py}, with certificate kind
\texttt{general\_integral\_morita\_rotation}. The full finite domains,
congruence vectors, separating vectors, and optimizer-free replay reports
are stored in \filename{research/round3/higher/}.
